\subsection{Code Release}
Our baseline code is attached in the submission as a minimal viable program. Full code will be released upon paper publication under the MIT license. 
This platform will be actively maintained and updated to support more features to support more research questions. We welcome any collaboration, contribution, feedback, and feature requests.

\subsection{Potential Risks}
Despite our efforts toward realism, our simulation operates in a constrained environment that inevitably omits many real-world factors. Results should therefore be treated as generating insights and hypotheses rather than definitive conclusions, and should not be used to directly guide important cognitive or policy decisions. Additionally, findings about conditions under which selfish strategies occasionally persist could, in principle, be interpreted as guidance for when self-interested behaviour might go unchecked---though the high level of abstraction in our simulation limits such applicability.

\subsection{More Discussion Over Our Methodology}
As we have emphasized, our method should be viewed as a complement to traditional mathematical models, not a replacement. By incorporating rich psychological realism into the simulation, our approach enables researchers to investigate how numerous factors interact in complex ways. However, this increased realism also means that simulation results are sensitive to the specific details of these factors and may not yield the definitive answers that highly abstract mathematical models can provide.

Yet, definitive answers are not always the primary goal of research, especially in the social sciences. Often, the objective is to uncover previously unnoticed factors that influence a phenomenon or to explore the intricate interplay among multiple variables. Such goals are difficult to achieve with traditional mathematical models, which require all relevant factors to be known or assumed in advance. Historically, researchers have relied on field studies to observe human behavior and identify these factors, but simulation now offers a cost-effective means to assist in discovery and hypothesis generation, potentially accelerating progress in the field.

Moreover, when the number of interacting factors becomes too great for analytical calculation, simulation becomes indispensable. While simulations inevitably deviate from reality—just as any modeling method, and such deviations may be amplified in large-scale runs—they can still provide valuable insights into research questions. Simulations can reveal what is possible, and the underlying mechanisms and developmental dynamics they expose may remain relevant even if the precise outcomes differ from those observed in the real world.

\subsection{Flexibility of the Simulation Platform}
Our platform is designed to be flexible and extensible. By varying the configuration settings, one can use the same platform to study different research questions. For example, we have used the same platform to study the effect of different moral types, different resource distributions, different communication costs, etc. In the section below, we also list a list of findings that are connected to different research areas that could possibly be investigated further with our platform.

Moreover, we support researchers to extend beyond morally related value dispositions. One can flexibly define the value dispositions of the agents by writing appropriate prompt templates. For example, one can define agents to be of different cultural backgrounds, different religions, different political views, etc. Or one can also study the effect of specific social norms by prescribing the agents to follow certain social rules, e.g, always equal distribution VS always contribution-based distribution, etc. Hunting-gathering environment equipped with general social interaction dynamics is very general to support a wide range of research questions.